%\documentclass[a4paper,addpoints,14pt]{exam}
\documentclass[answers,addpoints,14pt]{exam}
\unframedsolutions



%entorno para química: texto y math mode
%\newcommand*\chem[1]{\textup{#1}}
\newcommand*\chem[1]{\ensuremath{\mathrm{#1}}}

%tamaños fuentes extra: 8,9,14,17,20
	\usepackage{extsizes}

%display style in long fractions
\newcommand\ddfrac[2]{\frac{\displaystyle #1}{\displaystyle #2}}


%tabla generada
\usepackage[table,xcdraw]{xcolor}

%para los entornos de soluciones
\usepackage{mdframed,xcolor,color}
	%uso: 	entorno (begin,end): {mdframed}[backgroundcolor=brown!10] 


%lewis y química
\usepackage{chemformula,elements}
%chemformula, provides \ch{} with the !(<below>)(<formula>) syntax and \chlewis;
%elements, provides \elconf and \writeelconf.
%ex:
%\ch{
%  !(\elconf{Li})( "\chlewis{180.}{Li}" ) +
%  !(\elconf{F})( "\chlewis{0.90:180:270:}{F}" )
%   ->
%  !(\writeelconf{2})( Li+ ) +
%  !(\writeelconf{2,2+6})( "\chlewis{0:90:180:270:}{F}" {}- )
%}


%multicol
\usepackage{multicol}
  
  %símbolo bonito para qed
\usepackage{amsthm}
	%\newcommand{\mc}[3]{\multicolumn{#1}{#2}{#3}}
\newcommand{\halmos}{\vspace*{-4mm}
	
					\begin{flushright}
						\rule{1mm}{2.5mm} 
					\end{flushright}\vspace*{-4mm}
						
					  }
\renewcommand{\qed}{\halmos}

 
%página blanco
\usepackage{afterpage}
%uso \afterpage{\blankpage}

\newcommand\blankpage{%
    \null
    \thispagestyle{empty}%
    \addtocounter{page}{0}%
    \newpage}





 %caja bonita
  \usepackage[many]{tcolorbox}
\definecolor{titlebg}{RGB}{80,100,72}
\newtcolorbox{learn}{
  breakable,
  enhanced,
  arc=0pt,
  outer arc=0pt,
  colframe=titlebg,
  colback=titlebg!7,
  overlay unbroken and first={
    \node[
      draw=titlebg,
      fill=titlebg,
      rotate=270,
      anchor=north west,
      text=white,
      font=\bfseries
    ]
    at (frame.north west)  
    {TEST};
  }
}


\newtcolorbox{formulacion}{
  breakable,
  enhanced,
  arc=0pt,
  outer arc=0pt,
  colframe=titlebg,
  colback=titlebg!7,
  overlay unbroken and first={
    \node[
      draw=titlebg,
      fill=titlebg,
      rotate=270,
      anchor=north west,
      text=white,
      font=
    ]
    at (frame.north west)  
    {\hspace*{3ex}};
  }
}

\newtcolorbox{problemas}{
  breakable,
  enhanced,
  arc=0pt,
  outer arc=0pt,
  colframe=titlebg,
  colback=titlebg!7,
  overlay unbroken and first={
    \node[
      draw=titlebg,
      fill=titlebg,
      rotate=270,
      anchor=north west,
      text=white,
      font=
    ]
    at (frame.north west)  
    {PROBLEMAS};
  }
}
	
\usepackage{anysize}
	%\marginsize{left}{right}{top}{bottom}
    \marginsize{1cm}{.8cm}{1cm}{1cm}
    
\usepackage[spanish]{babel}
\usepackage[utf8]{inputenc}

\usepackage{sidecap}
\usepackage{mwe,float,afterpage}

%redefiniciones
\pointpoints{punto}{puntos}
\bonuspointpoints{punto}{puntos}
\hqword{Pregunta} %substitutes text for \Question"
%\vpgword{Pág.}%substitutestext for \Page"
\hpword{Máx.}%substitutes text for \Points"
\hsword{Nota}%substitutes text for \Score"
\htword{Total}%substitutes text for \Total:"

\chqword{Pregunta}
\chpword{Máx.}
\chsword{Nota}
\chtword{Total}
\chbpword{Puntos \textit{bonus}}

\renewcommand{\solutiontitle}{\noindent\textbf{Solución:}\par\noindent}


%fuentes
	%\usepackage[T1]{fontenc}
	%\usepackage{fourier}

	\usepackage[T1]{fontenc}
	\usepackage{librebaskerville}


\usepackage{amssymb,amsmath}
\usepackage{graphicx,graphics}
\usepackage{color}
	%	\colorgrids %grids con color
	%	\definecolor{GridColor}{gray}{0.8}
		%grid
%			\setlength{\gridsize}{5mm}
%   			 \setlength{\gridlinewidth}{0.1pt}

%fracciones en texto más elaboradas
	\usepackage{xfrac}

%mejora en la exposición de guarismos
	\usepackage{numprint}
 
%caption
\usepackage[font=small,labelfont=bf,
labelsep=endash]{caption}
	
%soluciones sombreadas	
\shadedsolutions


\checkboxchar{$\Box$}
\checkedchar{$\blacksquare$}

\usepackage{everypage,draftwatermark} %Marca de agua
\SetWatermarkText{\hspace*{-8cm} \parbox{\textwidth}{\centering  B1\\CIDEAD}}
\SetWatermarkScale{1.2}
\SetWatermarkLightness{0.9}
%
%\usepackage{background}
%\backgroundsetup{
%scale=1.5,
%angle=180,
%opacity=.05,
%color=black,
%contents={
%\includegraphics[scale=1]{unicorn.jpg}}}


\pagestyle{headandfoot}
%\runningheadrule
%\firstpageheader{Física para Biólogos}{Boletín 1a}{1 de octubre de 2011}
%\runningheader{Física para Biólogos}
%              {Boletín 1a, Página \thepage\ de \numpages}
%            {1 de octubre de 2011}
%\firstpagefooter{}{}{}
%\runningfooter{}{}{}


%espacios extra
%\extrawidth{1in}
%%\extrafootheight{.75in}
%\extraheadheight{1in}

%cabecera
\header{\includegraphics[scale=.22]{a.png}}       
      {}
       {$f_1$\textbf{B1-{\small CIDEAD}}\\\textbf{noviembre}2022}       
%pie de página

   %enumera en todas menos en la primera
   %\runningfootrule

   %enumera en todas
   %\footrule
   %\footer{}{Página \thepage\ de \numpages}{}
   \footer{}
	{Página \thepage\ de \numpages}
	{\iflastpage{{\footnotesize \textit{Fin del examen}}}{{\footnotesize \textit{Da la vuelta a la hoja}\ldots}}}


		
\begin{document}


%portada


%	\vspace{.5in}
%		
%	\makebox[\textwidth]{\textbf{Nombre y apellidos:}\enspace\hrulefill}
%	
%	\vspace{0.5in}
%
%%
%	\begin{center}
%		%\pointtable[h][questions]
%		\gradetable[h][questions]
%	\end{center}
%	
%	
%	\begin{center}
%		%\pointtable[h][questions]
%		%\combinedgradetable[h][questions]
%	\end{center}
%	
%		
%		
%	\vspace{0.3in}

%formato de preguntas
\qformat{\textbf{Pregunta \thequestion}\quad (\thepoints)\hfill} 


%\newpage
\vspace*{-2ex}

\fbox{\parbox{0.945\linewidth}{\textbf{Nombre y apellidos}:\vspace*{2mm}}}
\vspace{3ex}

%\textbf{\hspace*{5mm} Sólo una de las respuestas es correcta.\newline \hspace*{11mm} Las preguntas erróneamente contestadas restarán 0.25 puntos.}

\begin{learn}
{Cada respuesta incorrecta tendrá una penalización de $0.1$ puntos de examen. \textbf{¡No contestes a voleo!}}
\end{learn}	
		\vspace{0ex}
		
	\begin{questions}
	
%			
%	\begin{question}[1] ¿Cuál de las siguientes afirmaciones es \textbf{\small{INCORRECTA}}?
%		\begin{checkboxes} 
%			\choice $Z=18$ es un gas noble.
%			\CorrectChoice $Z=5$ tiene un electrón en su última capa.
%			\choice $Z=19$ es un metal.
%			\choice La configuración $1s^2\,2s^2p^6\,3s^2p^5$ pertenece a un elemento no metálico con siete electrones en su capa de valencia.
%		\end{checkboxes}
%		\end{question}
%		\qed
		
	
	\begin{question}[1] ¿Cuál de las siguientes afirmaciones es \textbf{\small{CORRECTA}}?
		\begin{checkboxes}
			\choice  Los metales en estado líquido no conducen la electricidad	.
			\CorrectChoice Las propiedades conductoras de los metales se explican por la presencia de una nube de electrones prácticamente libres.
			\choice La mayor parte de los metales son líquidos a temperatura ambiente debido a la nube electrónica.
			\choice Los metales sólidos tienen una estructura microscópica desordenada debido a la presencia de una nube de electrones.
		\end{checkboxes}
	\end{question}
	\qed
	\vspace{-2ex}
	
		
%	\begin{question}[1] Considera el compuesto al que da lugar el enlace entre los elementos $Z=20$, y $Z=9$. ¿Cuál de las siguientes afirmaciones es \textbf{\small{CORRECTA}}?
%		\begin{checkboxes} 
%		\choice Conduce la electricidad en estado sólido.
%		\choice Tiene un punto de fusión bajo.
%		\choice El enlace es covalente y forma una verdadera molécula.
%		\CorrectChoice Las intensas interacciones interatómicas provocarán que sea sólido a temperatura ambiente.
%		\end{checkboxes}
%		\end{question}
%		\qed
%	\vspace{-2ex}

	
%\begin{question}[1] Sean los elementos $_{17}X$, $_{20}Y$. ¿Qué afirmación es \textbf{\small{INCORRECTA}}?
%		\begin{checkboxes} 
%		\choice	Formarán un compuesto de fórmula \chem{YX_2}.
%		\choice \chem{Y} es un metal del grupo 2.
%		\choice \chem{YX_2} conduce la electricidad en disolución acuosa.
%		\CorrectChoice Formarán un compuesto de fórmula \chem{Y_2X}.
%		\end{checkboxes}
%		\qed
%		\end{question}
%		
	
	
%	\begin{question}[1] ¿Cuál de las siguientes afirmaciones es \textbf{\small{CORRECTA}} respecto al comportamiento del oxígeno, $_{\,8}\chem{O}$?
%		\begin{checkboxes}
%			\choice La existencia del oxígeno molecular puede justificarse por la compartición de un par de \textit{Lewis}.
%			\choice La molécula de $\chem{O_2}$ puede justificarse con un enlace iónico.
%			\choice Debido a que el oxígeno molecular es un gas, su existencia se justifica a través de la compartición masiva de electrones.
%			\CorrectChoice La existencia de la molécula de $\chem{O_2}$ puede justificarse a través de un enlace covalente con dos pares de \textit{Lewis} compartidos.
%		\end{checkboxes}
%	\end{question}
%	\qed
		
%
%	\begin{question}[1] ¿Cuál de las siguientes afirmaciones es \textbf{\small{VERDADERA}}?
%		\begin{checkboxes} 
%			\choice El enlace iónico une siempre un átomo no metálico con otro elemento no metálico.
%			\CorrectChoice El sulfato de amonio conduce la electricidad disuelto en agua.
%			\choice En una reacción química la suma del número de moles de reactivos es igual a la suma del número de moles de productos.
%			\choice Los compuestos moleculares son buenos conductores de la electricidad en estado líquido.
%		\end{checkboxes}
%		\qed
%		\end{question}
%	
%	
%\begin{question}[1] ¿Qué afirmación es {\small \textbf{INCORRECTA}} para el \textit{ácido sulfúrico}?
%	\begin{checkboxes} 
%		\choice Da lugar a sales ácidas al disociarse parcialmente.
%		\CorrectChoice Su molécula está formada por ocho átomos.
%		\choice Es uno de los tres ácidos ternarios que forma el azufre.
%		\choice El azufre actúa en él con estado de oxidación $6+$.
%	\end{checkboxes}
%\qed
%\end{question}

	
	\begin{question}[1] ¿Cuál de las siguientes afirmaciones es \textbf{\small{CORRECTA}} respecto al comportamiento del nitrógeno, $_{7}\chem{N}$?
		\begin{checkboxes}
			\choice La existencia del nitrógeno molecular puede justificarse por la compartición de un par de \textit{Lewis}.
			\choice La molécula de $\chem{N_2}$ puede justificarse con un enlace iónico.
			\choice Debido a que el nitrógeno molecular es un gas a temperatura ambiente, su existencia se justifica a través de la compartición masiva de electrones.
			\CorrectChoice La existencia de la molécula de $\chem{N_2}$ puede justificarse a través de un enlace covalente con tres pares de \textit{Lewis} compartidos.
		\end{checkboxes}
		
	\end{question}
	\qed
\vspace{-1ex}
		
%
%\question[1] ¿Cuál de las siguientes afirmaciones es {\small \textbf{INCORRECTA}}?
%		\begin{checkboxes} 
%		\choice El cloruro de amonio es sólido a temperatura ambiente.
%		\choice El comportamiento químico de un elemento viene determinado por los electrones de su última capa.
%		\choice El fluoruro de calcio es soluble en agua.
%		\CorrectChoice El tetracloruro de carbono es un buen conductor eléctrico. 
%		\end{checkboxes}
%\qed
%
%\vspace{-1ex}





\begin{question}[1] ¿Cuál de las siguientes afirmaciones es \textbf{\small{VERDADERA}}?
	\begin{checkboxes} 
			\choice El enlace covalente sólo se da entre metales.			
			\choice Todos los compuestos covalentes se encuentran en estado sólido a temperatura ambiente.	
			\CorrectChoice Los compuestos covalentes pueden formar redes cristalinas con puntos de ebullición muy altos.
			\choice Los sólidos covalentes conducen la electricidad disueltos.	
	\end{checkboxes}
\end{question}
\qed
	


%
%
%\question[3] Formula las siguientes sustancias	
%% Please add the following required packages to your document preamble:
%% \usepackage[table,xcdraw]{xcolor}
%% If you use beamer only pass "xcolor=table" option, i.e. \documentclass[xcolor=table]{beamer}
%\begin{table}[H]
%\centering
%\setlength{\tabcolsep}{12pt}
%\renewcommand{\arraystretch}{3}
%\begin{tabular}{|l|l|l|l|}
%\hline
%\rowcolor[HTML]{EFEFEF} 
%\multicolumn{1}{|c|}{\cellcolor[HTML]{EFEFEF}\textbf{Nombre}} & \textbf{Fórmula} & \multicolumn{1}{c|}{\cellcolor[HTML]{FFFFC7}\textbf{Nombre}} & \multicolumn{1}{c|}{\cellcolor[HTML]{FFFFC7}\textbf{Fórmula}} \\ \hline%Ácido peryódico                                               &                  & Dicromato de plata                                           &                                                               \\ \hline
%Ácido nítrico                                                 &                  & Ácido carbónico                                              &                                                               \\ \hline
%Nitrito de bario                                              &                  & Permanganato de amonio                                       &                                                               \\ \hline
%\end{tabular}
%\end{table}\qed
%
%\question[3] Nombra las siguientes especies químicas de las dos maneras propuestas.
%
%% Please add the following required packages to your document preamble:
%% \usepackage[table,xcdraw]{xcolor}
%% If you use beamer only pass "xcolor=table" option, i.e. \documentclass[xcolor=table]{beamer}
%\begin{table}[H]
%\renewcommand{\arraystretch}{3}
%\begin{tabular}{|l|l|l|}
%\hline
%\rowcolor[HTML]{EFEFEF} 
%\multicolumn{1}{|c|}{\cellcolor[HTML]{EFEFEF}\textbf{Compuesto}} & \multicolumn{1}{c|}{\cellcolor[HTML]{EFEFEF}\textbf{Nomenclatura tradicional}} & \multicolumn{1}{c|}{\cellcolor[HTML]{EFEFEF}\textbf{Nomenclatura de hidrógeno}} \\ \hline
%
%$\chem{Ni(NO_2)_3}$                                              &                                                                                &                                                                                 \\ \hline
%$\chem{Ag_2SO_3}$                                                &                                                                                &                                                                                 \\ \hline
%$\chem{NaIO_4}$                                                  &                                                                                &                                                                                 \\ \hline
%
%$\chem{NH_4(HSO_4)}$                                             &                                                                                &                                                                                 \\ \hline
%$\chem{HBrO}$                                                    &                                                                                &                                                                                 \\ \hline
%$\chem{FeSO_4}$                                                 &                                                                                &                                                                                 \\ \hline
%\end{tabular}
%\end{table}
%
%\qed
%
%
%
%\newpage

\vspace*{3ex}



\begin{formulacion}
\textbf{FORMULACIÓN Y NOMENCLATURA}
\end{formulacion}


\begin{question}[4]
	Nombra cada una de estas sustancias \textbf{de dos formas}, siempre que no coincidan. 
	\begin{itemize}
		\item $\chem{HMnO_4}$
		%\vspace{50pt}
		\begin{mdframed}[backgroundcolor=brown!10]  
		Ácido permangánico/ Tetraoxidomanganato de hidrogeno
		\end{mdframed}
		
%		\item $\chem{HClO_2}$
%		\vspace{60pt}
				
		\item $\chem{FeSO_4}$
		%\vspace{50pt}
		\begin{mdframed}[backgroundcolor=brown!10]  
		Sulfato de hierro(II)/ Tetraoxidosulfato de hierro
		\end{mdframed}
		%\item $\chem{HIO}$\vspace{50pt}
		
		\item $\chem{CaCO_3}$
		%\vspace{50pt}
		\begin{mdframed}[backgroundcolor=brown!10]  
		Carbonato de calcio/ Trioxidocarbonato de calcio
		\end{mdframed}
		\item $\chem{HNO_3}$
		%\vspace{50pt}
		\begin{mdframed}[backgroundcolor=brown!10]  
		Ácido nítrico/ Trioxidonitrato de hidrogeno
		\end{mdframed}
%		\item $\chem{HCl(aq)}$
%		\vspace{60pt}
	
	\end{itemize}
\end{question}



\begin{question}[4] Formula las siguientes sustancias. 

\begin{mdframed}[backgroundcolor=brown!10]  
Si tienes dudas sobre ello, consúltame.
\end{mdframed}
		\begin{itemize}
		\item 1,3,5-tribromobenceno%\vspace{50pt}
		
		\item Butanamida%\vspace{50pt}
		
		\item Propanodial%\vspace{50pt}
		
		\item Etanoato de metilo%\vspace{50pt}
		
		\item 2-fenilfenol%\vspace{50pt}
				
		\item  Ácido pent-2-en-dioico%\vspace{50pt}
		
		%\item  But-3-in-1,2-diol\vspace{50pt}
		
	%	\item  Ciclopentilciclohexano \vspace{70pt}
		
	\end{itemize}
\end{question}
\newpage


	
\begin{problemas}{Recuerda que un examen es una \textit{comunicación}: debes invocar las leyes en que te basas para resolver los ejercicios así como \textit{indicar brevemente los pasos y razonamientos} que sigues para obtener los resultados.}
\end{problemas}
\vspace{5ex}


\question[4] 
	\begin{parts}
	  \setlength\itemsep{1ex}
		\part[2] Representa la \textbf{estructura de niveles y subniveles} del átomo de hidrógeno para $n=1$, $n=2$, $n=3$ y $n=4$. 
		\begin{mdframed}[backgroundcolor=brown!10]  
		Se realiza en un diagrama de tipo árbol sin más que saber que \textit{para cada} valor permitido de $l$ hay que asociar una serie de posibles valores de $m$. A modo de ejemplo, para $n=4$ son posibles los valores $l\in\{0,1,2,3\}$ y \textit{para cada uno} de estos valores, $m\in\{-l,\dots,+l\}$.
		
		{
		\centering
		\begin{tabular}{|c|c|c|}
		\hline 
		$n=4$ &  &  \\ 
		\hline 
		 & $l=3$ & $m\in\{-3,-2,-1,0,+1,+2,+3\} $\\ 
		\hline 
		 & $l=2$ & $m\in\{-2,-1,0,+1,+2\}$ \\ 
		\hline 
		 & $l=1$ & $m\in\{-1,0,+1\} $\\ 
		\hline 
		 & $l=0 $& $m=0$ \\ 
		\hline 
		\end{tabular}\par
							}
		\end{mdframed}
		\part[2] Calcula razonadamente el \textbf{número máximo de electrones} que podría dar cabida cada uno de los niveles anteriores.
		\begin{mdframed}[backgroundcolor=brown!10]  
		Teniendo en cuenta el principio de exclusión de Pauli en cada uno de los orbitales descritos por la terna $(n,m,l$ cabrán exactamente dos electrones, correspondientes a los dos posibles valores del número cuántico $m_s=\pm\sfrac{1}{2}$. 
		
		El número de electrones que caben en $n=1$ será, por tanto, $2$; para $n=2$ tendremos $8\,e^-$; en $n=3$, caben $18\;e^-$ y en $n=4$, resulta que pueden alojarse $32\;e^-$. 
		
		En general puedes intentar demostrar (no es difícil: solo hay que sumar una serie aritmética) que el número de electrones en un nivel etiquetado por $n$ es de $2n^2$. 
		\end{mdframed}
	\end{parts}
\qed
	
	\newpage
	
\question[6] Considera los elementos desconocidos $^{8}\alpha$, $^{18}\beta$, $^{34}\delta$ y $^{38}\epsilon.$ 
    	\begin{parts}
    	\setlength\itemsep{1ex}
    		\part[1] Escribe sus \textbf{configuraciones electrónicas}. 
    			\begin{mdframed}[backgroundcolor=brown!10]  
    				\begin{itemize}
    					\item $[^8\alpha]=\elconf{O}$
    					\item $[^{18}\beta]=\elconf{Ar}$
    					\item $[^{34}\delta]=\elconf{Se}$
    					\item $[^{38}\epsilon]=\elconf{Sr}$
    				\end{itemize}
    			\end{mdframed}
    		\part[2] Justifica el \textbf{tipo de enlace y la fórmula del compuesto} que forman los elementos $\delta$ y $\epsilon$.
    		\begin{mdframed}[backgroundcolor=brown!10]  
    		A la vista de la configuración $\delta$ es un no metal del grupo $16$ que \textit{busca} ganar dos electrones para cerrar su capa de valencia.   		
    		 Por otro lado $\epsilon$ es un metal alcalinotérreo, electropositivo, y que \textit{busca} perder dos electrones para encontrar la estabilidad. 	Se unirán por enlace iónico dando lugar al compuesto $\epsilon\delta$.
    		\end{mdframed}
    		\part[1] \textbf{Describe las propiedades} (estructura, solubilidad, conducción eléctrica y estado físico más probable a temperatura ambiente) que presentará el compuesto del apartado b).
    		\begin{mdframed}[backgroundcolor=brown!10]  
    		Como compuesto iónico probablemente será un sólido a temperatura ambiente debido a la fuerte interacción de carácter electromagnético que actúa interatómicamente en la red cristalina. Será soluble en disolventes polares y conducirá la electricidad disuelto o fundido, estados en que dispondrá de cargas móviles. 
    		\end{mdframed}
    		\part[1] \textbf{Justifica la existencia} de la molécula $\alpha_2$ e indica el \textbf{tipo de enlace} que se produce.
    		\begin{mdframed}[backgroundcolor=brown!10]  
    		$\alpha$ es un no metal del grupo $16$. Para conseguir la estabilidad puede interaccionar con un átomo de la misma especie compartiendo un par de electrones (par de \textit{Lewis}) de modo que consiga alcanzar una estructura electrónica de gas noble.
    		\end{mdframed}
    		\part[1] Justifica el \textbf{tipo de enlace} que formarán $\beta$ y $\alpha$.
    		\begin{mdframed}[backgroundcolor=brown!10]  
    		$\beta$ es un gas noble y no forma compuestos.
    		\end{mdframed}
    	\end{parts}
		\qed



%\afterpage{\blankpage}
\newpage
\vspace*{-7ex}

\begin{center}
       \fbox{\footnotesize HOJA DE RESPUESTAS A LOS PROBLEMAS}
\end{center}

\end{questions}	


\end{document}






